%\subsection[\texttt{python} side]{\python side}\label{sec:python_side}
All binary bindings have their counterparts
on the \python side, with access provided by
two \python scripts: \texttt{gretl4py\_addons.py}
and \texttt{gretl4py\_classes.py}. Additionally,
\gretlpy is initialized by the \texttt{\_\_init\_\_.py} file.


\subsection{Module related tools}\label{sec:module_related_python}
\gretlfpycmd{GRETL4PY\_VERBOSE}\newline
\generaldesc{An environmental variable that controls
	the printing of debug information generated by \gretlpy.}

\gretlfpycmd{version}
\estdesc{The \gretlpy module version.}
\begin{lstlisting}[style=pythonstylefoot]
gretl.version -> str
\end{lstlisting}

\gretlfpycmd{about()}
\estdesc{This function retrieves and optionally prints
	basic information about the \gretlpy module, including
	the libraries it utilizes and the contents of \lgretl's
	\gretlcmd{\$sysinfo} bundle.}
\begin{lstlisting}[style=pythonstylefoot]
gretl.about(do_print: bool = False) -> gretl._gretl.Bundle
\end{lstlisting}
\generaldesc{The returned dictionary contains the following
	entries (alphabetical order):
	\begin{itemize}
		\item \pythcode{built\_in\_strings}: The values of
			\lgretl's built-in strings \citep[see][for details]{gretl_manual}
			(\pythcode{\_gretl.Bundle}).
		\item \pythcode{lg\_ver}: The version of the \lgretl[-1.0]
			library (\pythcode{str}).
		\item \pythcode{mod\_ver}: The version of the \gretlpy
			module (\pythcode{str}).
		\item \pythcode{py11\_ver}: The version of the \pkg{pybind11}
			module, formatted as \texttt{MAJOR.MINOR.PATCH} (\pythcode{str}).
		\item \pythcode{sysinfo}: The contents of \lgretl's
			\gretlcmd{\$sysinfo} bundle (\pythcode{\_gretl.Bundle}).
		\end{itemize}}


\subsection{The \gretlpy functions}\label{sec:gretl4py_functions}
\vspace*{10pt}\noindent The \gretlpy module provides the following functions:

\gretlfpycmd{cdf()}\label{func:gretl.cdf}
\estdesc{Cumulative distribution function calculator.}
\begin{lstlisting}[style=pythonstylefoot]
gretl.cdf(d: str, x: float[, param: object = None]) -> float
gretl.cdf(d: str, x: array_like[, param: object = None]) -> numpy.ndarray
gretl.cdf(d: str, x: str[, param: object = None]) -> numpy.ndarray
\end{lstlisting}
\maptof{cdf(d, x)}{float | numpy.ndarray}
\maptof{cdf(d, param, x)}{float | numpy.ndarray}
Examples:\newline
\pythcodef{p1 = gretl.cdf('N', -2.5)}\\
\pythcodef{p2 = gretl.cdf('X', 5.67, 3)}\\
\pythcodef{p3 = gretl.cdf('D', 0.25, [-1, 1])}\\
\generaldesc{Returns $P(X \leq x)$, where the distribution
	of \textit{X} is determined by the string \textit{d}.}\\
If \textit{x} is of type:
\begin{itemize}
	\item \pythcode{list}, \pythcode{tuple}, or \pythcode{NumPy array},
	it will be converted internally to \texttt{numpy.ndarray}
	of type \texttt{float64};
	\item \pythcode{str}, it is interpreted as a \gretlcmd{series} name.
	Note that \pythcode{gretl.cdf()} works by default only with the
	default dataset. To obtain CDFs for variables from different
	datasets, please refer to \ref{met:gretldataset::cdf}.
\end{itemize}
The \textit{param} argument consists of zero or more additional scalar
arguments required to specify the parameters of the distribution, as
follows:
\begin{footnotesize}\begin{itemize}
	\item \textbf{Standard normal} (\texttt{d = z, n, N}): no additional arguments
	\item \textbf{Bivariate normal} (\texttt{D}): correlation coefficient
	\item \textbf{Logistic} (\texttt{lgt} or \texttt{s}): no additional arguments
	\item \textbf{Student's $t$} (\texttt{t}): degrees of freedom
	\item \textbf{Chi-square} (\texttt{c}, \texttt{x}, or \texttt{X}): degrees of freedom
	\item \textbf{Snedecor's $F$} (\texttt{f} or \texttt{F}): degrees of freedom (numerator); degrees of freedom (denominator)
	\item \textbf{Gamma} (\texttt{g} or \texttt{G}): shape; scale
	\item \textbf{Beta} (\texttt{beta}): two shape parameters
	\item \textbf{Binomial} (\texttt{b} or \texttt{B}): probability; number of trials
	\item \textbf{Poisson} (\texttt{p} or \texttt{P}): mean
	\item \textbf{Exponential} (\texttt{exp}): scale
	\item \textbf{Weibull} (\texttt{w} or \texttt{W}): shape; scale
	\item \textbf{Laplace} (\texttt{l} or \texttt{L}): mean; scale
	\item \textbf{Generalized Error} (\texttt{E}): shape
	\item \textbf{Non-central chi-square} (\texttt{ncX}): degrees of freedom; non-centrality parameter
	\item \textbf{Non-central $F$} (\texttt{ncF}): degrees of freedom (numerator); degrees of freedom (denominator); non-centrality parameter
	\item \textbf{Non-central $t$} (\texttt{nct}): degrees of freedom; non-centrality parameter
\end{itemize}\end{footnotesize}

\gretlfpycmd{critical()}\label{func:gretl.critical}
\estdesc{Critical value calculator.}
\begin{lstlisting}[style=pythonstylefoot]
gretl.critical(c: str, p: float[, param: object = None]) -> float
gretl.critical(c: str, p: array_like[, param: object = None]) -> numpy.ndarray
gretl.critical(c: str, p: str[, param: object = None]) -> numpy.ndarray
\end{lstlisting}
\maptof{critical(c, p)}{float | numpy.ndarray}
\maptof{critical(c, param, p)}{float | numpy.ndarray}
Examples:\newline
\pythcodef{c1 = gretl.critical('t', 0.025, 20)}\\
\pythcodef{c2 = gretl.critical('F', 0.05, [4, 48])}\\
\generaldesc{Returns \textit{x} such that $P(X > x) = p$, where
	the distribution \textit{X} is determined by the character \textit{c}.}
For details about the meaning of the arguments, please
refer to \ref{func:gretl.cdf}. Note that \pythcode{gretl.critical()}
works by default only with the default dataset. To obtain critical
values for variables from different datasets, please refer to
\ref{met:gretldataset::critical}.

\gretlfpycmd{delete()}
\estdesc{This is an all-purpose destructor. It should be used with caution; no confirmation is asked.}
\begin{lstlisting}[style=pythonstylefoot]
gretl._gretl.delete(varname: str)
\end{lstlisting}
\mapto[clean]{delete varlist}
\mapto[clean]{delete varname}
\mapto[]{delete --type=type-name}
\begin{itemize}
	\item[]\maptoopt{type=str}{--type=type-name}{\seegretlmanual}
	\item[]\maptoopt{db=True}{db}{delete series from database}
	\item[]\pythcode{data=gretl.Dataset}\par%
		\hspace*{10pt}{\footnotesize\pythcode{gretl.delete()} works with the specified dataset}
\end{itemize}
Examples:\newline
\pythcodef{.delete(varlist)}\\
\pythcodef{.delete(varname)}\\
\pythcodef{.delete(type='type-name')}\\
\generaldesc{In the first form above, varlist is a list of series,
	given by name or ID number. Note that when you delete series
	any series with higher ID numbers than those on the deletion
	list will be re-numbered.\newline
	In the second form, the name of a \texttt{scalar}, \texttt{matrix},
	\texttt{string} or \texttt{bundle} may be given for deletion.\newline
	In the third form, the \gretlcmd{--type} option must be accompanied
	 by one of the following type-names: \texttt{matrix}, \texttt{bundle},
	 \texttt{string}, \texttt{list}, \texttt{scalar} or \texttt{array}.}

\gretlfpycmd{df2dict()}
\begin{lstlisting}[style=pythonstylefoot]
gretl.df2dict(dframe: pandas.core.frame.DataFrame[, pd: int, with_obs: bool = False]) -> dict
\end{lstlisting}
\generaldesc{This function converts a \pkg{pandas}
	\texttt{DataFrame} file into a dictionary compatible with
	\gretlpy (as described in Section~\ref{sec:data_management_python})
	and accepts two optional arguments:
	\begin{itemize}
		\item \pythcode{pd}: If provided, it
		specifies the data's periodicity (default:~1).
		\item \pythcode{with\_obs}: If set to \pythcode{True},
		the first column of the given \texttt{DataFrame} is
		treated as containing observation labels).
\end{itemize}}

\gretlfpycmd{fc\_print()}
\estdesc{XXX}
\begin{lstlisting}[style=pythonstyle]
gretl.fc_print(fc=dict)
\end{lstlisting}

\gretlfpycmd{genr()}
\estdesc{NOTE: for comprehensive and updated info on this
	function you'll want to refer to \cite[chapter 10]{gretl_manual}.}
\begin{lstlisting}[style=pythonstylefoot]
gretl.genr(formula=str, data=gretl._gretl.Dataset)
gretl.genr(name=str, value=numpy.ndarray)
gretl.genr(name=str, value=int)
gretl.genr(name=str, value=float)
gretl.genr(name=str, value=str)
\end{lstlisting}
\mapto[clean]{genr newvar = formula}

\gretlfpycmd{get\_accessor()}
\estdesc{Function to read \lgretl['s] accessors}
\begin{lstlisting}[style=pythonstylefoot]`
gretl.get_accessor(name: str) -> object
\end{lstlisting}
Examples:\newline
\pythcodef{.get\_accessor('\$version')}\\
\pythcodef{.get\_accessor('\$stopwatch')}\\
\generaldesc{Note: \pythcode{gretl.get\_accessor()} works by default only
	with the default dataset. To access accessors related to different datasets
	or models, please refer to \ref{met:gretldataset::get_accessor}
	or \ref{met:gretlmodel::get_accessor}.}

\gretlfpycmd{get\_data()}
\estdesc{Loads provided data as \pythcode{gretl._gretl.Dataset}}
\begin{lstlisting}[style=pythonstylefoot]
gretl.get_data(src: str | dict) -> gretl._gretl.Dataset
\end{lstlisting}
\generaldesc{This function recognizes the following data
	formats:
	\begin{itemize}
		\item \gretl native formats: \texttt{gdt}, and \texttt{gdtb},
		\item text files: \texttt{csv},
		\item formats known from other econometric programs: \texttt{dta}, and \texttt{wf1},
		\item spreadsheets: \texttt{xls}, \texttt{xlsx}, and \texttt{ods}
		\item \python['s] \pythcode{dict} as defined in \ref{gretl4py:data:dict}.
	\end{itemize}
	For a complete list of supported please refer to \cite{gretl_manual}.\newline
	Note: the first loaded dataset is automatically registered
	as the global dataset (see \ref{sec:globalandindivodual}).}

\gretlfpycmd{include()}
\estdesc{Intended for including gretl function packages (see \ref{sec:classpypackage}).}
\begin{lstlisting}[style=pythonstylefoot]
gretl.include(filename: str) -> gretl._gretl.Package
\end{lstlisting}
\mapto[clean]{include filename}
\generaldesc{\textit{filename} should have the extension \textit{gfn}
	(a gretl function package) or no extension. Functions provided by
	\textit{filename} are bound in user space.}

\gretlfpycmd{invcdf()}\label{func:gretl.invcdf}
\estdesc{Inverse cumulative distribution function calculator.}
\begin{lstlisting}[style=pythonstylefoot]
gretl.invcdf(d: str, u: float[, param: object = None]) -> float
gretl.invcdf(d: str, u: array_like[, param: object = None]) -> numpy.ndarray
gretl.invcdf(d: str, u: str[, param: object = None]) -> numpy.ndarray
\end{lstlisting}
\maptof{invcdf(d, u)}{float | numpy.ndarray}
\maptof{invcdf(d, param, u)}{float | numpy.ndarray}
\generaldesc{For a continuous distribution, returns \textit{x}
	such that $P(X \leq x) = u$, for $u \in [0, 1]$. For a discrete
	distribution (Binomial or Poisson), returns the smallest
	\textit{x} such that $P(X \leq x) \geq u$.}\\
For details about the meaning of the arguments, please
refer to \ref{func:gretl.cdf}. Note that \pythcode{gretl.invcdf()}
works by default only with the default dataset. To obtain
inverted CDFs for variables from different datasets, please
refer to \ref{met:gretldataset::invcdf}.

\gretlfpycmd{modeltab()}
\estdesc{Prints a \gretl-like ``model table''.}
\begin{lstlisting}[style=pythonstylefoot]
gretl.modeltab(models: list[, output: str="print", tablefmt: str="github", stars: bool=True, coeff_fmt: str=".3f", stderr_fmt: str=".4f"])
\end{lstlisting}
\begin{footnotesize}
where options are:
\begin{itemize}
	\item[] \textit{models}: list of GretlModel objects,
	\item[] \textit{output}: \texttt{``print''}---print the table,
		\texttt{``return''}---return as string,
	\item[] \textit{tablefmt}: table format for \texttt{tabulate} package,
	\item[] \textit{stars}: whether to add significance stars,
	\item[] \textit{coeff\_fmt}: format string for coefficients,
	\item[] \textit{stderr\_fmt}: rormat string for standard errors.
\end{itemize}\end{footnotesize}

\gretlfpycmd{nparray2dict()}
\begin{lstlisting}[style=pythonstylefoot]
gretl.nparray2dict(X: numpy.ndarray[, pd: int = 1, Y: numpy.ndarray, with_const: bool = False, const_pos: int = 1]) -> dict
\end{lstlisting}
\generaldesc{This function converts a \pkg{numpy}
	array file into a dictionary compatible with
	\gretlpy (as described in Section~\ref{sec:data_management_python})
	and accepts four optional arguments:
	\begin{itemize}
		\item \pythcode{pd}: If provided, it
		specifies the data's periodicity.
		\item \pythcode{Y}: If provided, it
		contains observations for the dependent variable.
		\item \pythcode{with\_const}: If set to
		\pythcode{True}, it specifies that the provided
		array \textit{X} contains a constant term.
		\item \pythcode{const\_pos}: Used only if
		\pythcode{with\_const=True}, and specifies which.
\end{itemize}}

\gretlfpycmd{pdf()}\label{func:gretl.pdf}
\estdesc{Probability density function calculator.}
\begin{lstlisting}[style=pythonstylefoot]
gretl.pdf(d: str, x: float[, param: object = None]) -> float
gretl.pdf(d: str, x: array_like[, param: object = None]) -> numpy.ndarray
gretl.pdf(d: str, x: str[, param: object = None]) -> numpy.ndarray
\end{lstlisting}
\maptof{pdf(d, x)}{float | numpy.ndarray}
\maptof{pdf(d, param, x)}{float | numpy.ndarray}
Examples:\newline
\pythcodef{f1 = gretl.pdf('N', -2.5)}\\
\pythcodef{f2 = gretl.pdf('X', x, 3)}\\
\pythcodef{f3 = gretl.pdf('W', x, [shape, scale])}\\
\generaldesc{Returns the density at \textit{x} of the
	distribution identified by the code \textit{d}.}\\
For details about the meaning of the arguments, please
refer to \ref{func:gretl.cdf}. Note that \pythcode{gretl.pdf()}
works by default only with the default dataset. To obtain PDFs
for variables from different datasets, please refer to \ref{met:gretldataset::pdf}.

\gretlfpycmd{print()}
\estdesc{Please note that print is a rather ``basic'' function
	(primarily intended for printing the values of series). For
	more flexible printers please refer to \ref{met:gretlbundle::print},
	\ref{met:gretldataset::print}.}
\begin{lstlisting}[style=pythonstylefoot]
gretl.print(varlist: str[, OPTS])
\end{lstlisting}
\mapto[clean]{print varlist}
\mapto[]{print}
\begin{itemize}
	\item[]\maptoopt{byobs=True}{--byobs}{by observations}
	\item[]\maptoopt{no\_dates=True}{--no-dates}{use simple observation numbers}
	\item[]\maptoopt{range='start:stop'}{--range=start:stop}{\seegretlmanual}
	\item[]\maptoopt{head=lines}{--head[=lines]}{\seegretlmanual}
	\item[]\maptoopt{tail=True}{--tail[=lines]}{\seegretlmanual}
	\item[]\maptoopt{midas=True}{--midas}{\seegretlmanual}
	\item[]\pythcode{data=gretl.Dataset}\par%
		\hspace*{10pt}{\footnotesize\pythcode{gretl.print()} works with the specified dataset}
\end{itemize}
Examples:\newline
\pythcodef{.print('y x', byobs=True)}\\
\pythcodef{.print('my\_matrix')}\\
\pythcodef{.print('my\_array', range='3:6')}\\
\pythcodef{.print('hflist', midas=True)}

\generaldesc{In the first variant shown above (also see the
	first example), \textit{varlist} should be a list of series
	(either a named list or a list specified via the names
	or ID numbers of series, separated by spaces).\newline
	If no argument is given (the second variant shown above)
	then the action is similar to the first case except that
	all series in the current dataset are printed. The supported
	options are as described above. }

\gretlfpycmd{pvalue()}\label{func:gretl.pvalue}
\estdesc{\textit{P}-value calculator.}
\begin{lstlisting}[style=pythonstylefoot]
gretl.pvalue(c: str, x: float[, param: object = None]) -> float
gretl.pvalue(c: str, x: array_like[, param: object = None]) -> numpy.ndarray
gretl.pvalue(c: str, x: str[, param: object = None]) -> numpy.ndarray
\end{lstlisting}
\maptof{pvalue(c, x)}{float | numpy.ndarray}
\maptof{pvalue(c, param, x)}{float | numpy.ndarray}
Examples:\newline
\pythcodef{p1 = gretl.pvalue('z', 2.2)}\\
\pythcodef{p2 = gretl.pvalue('X', 5.67, 3)}\\
\pythcodef{p3 = gretl.pvalue('F', 5.67, [3, 30])}\\
\generaldesc{Returns $P(X > x)$, where the distribution of
	\textit{X} is determined by the character \textit{c}.}\\
For details about the meaning of the arguments, please
refer to \ref{func:gretl.cdf}. Note that \pythcode{gretl.pvalue()}
works by default only with the default dataset. To obtain \textit{P}-values
for variables from different datasets, please refer to \ref{met:gretldataset::pvalue}.

\gretlfpycmd{print\_module\_source()}
\estdesc{XXX}
\begin{lstlisting}[style=pythonstyle]
	gretl.print_module_source(object)
\end{lstlisting}

\gretlfpycmd{series2dict()}
\begin{lstlisting}[style=pythonstylefoot]
gretl.series2dict(Y : pandas.core.series.Series[, pd: int = 1, varname: str = 'y']) -> dict
\end{lstlisting}
\generaldesc{This function converts a \pkg{pandas}
	series into a dictionary compatible with
	\gretlpy (as described in Section~\ref{sec:data_management_python})
	and accepts two optional arguments:
	\begin{itemize}
		\item \pythcode{pd}: If provided, it
		specifies the data's periodicity.
		\item \pythcode{varname}: If provided, it
		contains the variable's name).
\end{itemize}}

\gretlfpycmd{set()}
\estdesc{The most common use of this function is the first
	map variant shown below, where it is used to set the value
	of a selected \lgretl parameter.}
\begin{lstlisting}[style=pythonstylefoot]
gretl.set()
gretl.set(variable: str = ''[, value: str = '', OPTS])
gretl.set(variable: str, value: int[, OPTS])
\end{lstlisting}
\mapto[0]{set variable value}
\mapto[0]{set --to-file=filename}
\mapto[0]{set --from-file=filename}
\mapto[0]{set stopwatch}
\mapto[]{set}
\begin{itemize}
	\item[]\pythcode{data=gretl.Dataset}\par%
		\hspace*{10pt}{\footnotesize\pythcode{gretl.set()} works with the specified dataset}
\end{itemize}
\generaldesc{The other uses are: with \gretlcmd{--to-file} to write
	a script file containing all the current parameter settings;
	with \gretlcmd{--from-file} to read a script file containing
	parameter settings and apply them to the current session; with
	stopwatch to zero the gretl ``stopwatch'' which can be used to
	measure CPU time (\seegretlmanual the entry for the
	\gretlcmd{\$stopwatch} accessor); or, if the word set is given alone,
	to print the current settings.}

\gretlfpycmd{show\_data\_status()}
\estdesc{List of available datasets.}
\begin{lstlisting}[style=pythonstylefoot]
gretl.show_data_status()
\end{lstlisting}

\gretlfpycmd{summary()}
\estdesc{In its first form, this command prints summary statistics
	for the variables in \textit{varlist}, or for all the variables
	in the data set if \textit{varlist} is omitted.}
\begin{lstlisting}[style=pythonstylefoot]
gretl.summary(varlist: str[, OPTS])
\end{lstlisting}
\mapto[0]{summary [ varlist ]}
\mapto{summary --matrix=matname}
\begin{itemize}
	\item[]\maptoopt{simple=True}{--simple}{basic statistics only}
	\item[]\maptoopt{weight='wtvar'}{--weight=wtvar}{weighting variable}
	\item[]\maptoopt{by='byvar'}{--by=byvar}{\seegretlmanual}
	\item[]\pythcode{data=gretl.Dataset}\par%
	\hspace*{10pt}{\footnotesize\pythcode{gretl.summary()} works with the specified dataset}
\end{itemize}
\generaldesc{The table of statistics produced by \gretlcmd{summary}
	can be retrieved in matrix form via the \gretlcmd{\$result}
	accessor. When the \gretlcmd{--by} option is given, this accessor
	is produced only if \textit{varlist} contains a single series.}

\gretlfpycmd{uservars\_list()}
\estdesc{XXX}
\begin{lstlisting}[style=pythonstyle]
	gretl._gretl.uservars_list(void)
\end{lstlisting}
\generaldesc{XXX}

